\documentclass{article}
\usepackage{amsmath,amssymb,amsthm}
\usepackage{algorithm}
\usepackage{algpseudocode}
\usepackage{enumitem}
\usepackage{hyperref}
\usepackage{bm}
\newtheorem{theorem}{Theorem}
\newtheorem{lemma}{Lemma}
\newtheorem{assumption}{Assumption}
\newcommand{\E}{\mathbb{E}}
\newcommand{\R}{\mathbb{R}}

\begin{document}

\section*{Algorithm Name}
\textbf{Probabilistic Momentum Coordinate Descent (PMCD)}  

The method is a coordinate‑wise variant of the heavy‑ball momentum scheme.  
At iteration $k$ a single coordinate $i_k\in\{1,\dots,d\}$ is drawn from a fixed
distribution $p=(p_1,\dots,p_d)$, $p_i>0$, $\sum_i p_i=1$.  
Only the selected coordinate is updated, but a momentum buffer is kept for
every coordinate.

---------------------------------------------------------------------

\section*{Mathematical Formulation}
Let $f:\R^d\rightarrow\R$ be differentiable.  
Denote by $\nabla_i f(w)$ the partial derivative w.r.t. coordinate $i$.
For each coordinate $i$ we maintain a momentum variable $v_i^k\in\R$.
Given step size $\eta>0$ and momentum parameter $\beta\in[0,1)$, the
iterations are

\[
\begin{aligned}
\text{Sample } i_k &\sim p,\\[4pt]
v_i^{k+1} &=
\begin{cases}
\beta\,v_i^{k}+ \nabla_i f(w^{k}), &\text{if } i=i_k,\\[4pt]
v_i^{k}, &\text{otherwise},
\end{cases}\\[6pt]
w_i^{k+1} &=
\begin{cases}
w_i^{k}-\eta\, v_i^{k+1}, &\text{if } i=i_k,\\[4pt]
w_i^{k}, &\text{otherwise}.
\end{cases}
\end{aligned}
\tag{PMCD}
\]

All coordinates not selected remain unchanged.

---------------------------------------------------------------------

\section*{Pseudocode}
\begin{algorithm}[H]
\caption{Probabilistic Momentum Coordinate Descent}
\begin{algorithmic}[1]
\State \textbf{Input:} initial point $w^{0}\in\R^{d}$, step size $\eta>0$, momentum $\beta\in[0,1)$,
probability vector $p\in\Delta_{d}$, maximum iterations $T$.
\State Initialize $v^{0}\gets 0\in\R^{d}$.
\For{$k=0$ \textbf{to} $T-1$}
    \State Sample $i_k\sim p$.
    \State $v_{i_k}^{k+1}\gets \beta\,v_{i_k}^{k}+ \nabla_{i_k}f(w^{k})$.
    \State $w_{i_k}^{k+1}\gets w_{i_k}^{k}-\eta\,v_{i_k}^{k+1}$.
    \For{all $i\neq i_k$}
        \State $v_i^{k+1}\gets v_i^{k}$,\quad $w_i^{k+1}\gets w_i^{k}$.
    \EndFor
\EndFor
\State \textbf{Output:} $w^{T}$.
\end{algorithmic}
\end{algorithm}

---------------------------------------------------------------------

\section*{Dry Run Example}
Consider the scalar quadratic $f(w)=(w-3)^{2}$, thus $d=1$,
$\nabla f(w)=2(w-3)$.  
Take $w^{0}=0$, step size $\eta=0.1$, momentum $\beta=0.5$, and $p_{1}=1$ (the only
coordinate is always chosen).

\begin{enumerate}[label=Iteration \arabic*:, leftmargin=*]
\item $k=0$: 
\[
\begin{aligned}
v^{1}&=\beta v^{0}+ \nabla f(w^{0})
      =0+2(0-3)=-6,\\
w^{1}&=w^{0}-\eta v^{1}=0-0.1(-6)=0.6,\\
f(w^{1})&=(0.6-3)^{2}=5.76.
\end{aligned}
\]

\item $k=1$: 
\[
\begin{aligned}
v^{2}&=\beta v^{1}+ \nabla f(w^{1})
      =0.5(-6)+2(0.6-3)= -3-4.8=-7.8,\\
w^{2}&=w^{1}-\eta v^{2}=0.6-0.1(-7.8)=1.38,\\
f(w^{2})&=(1.38-3)^{2}=2.6244.
\end{aligned}
\]

\item $k=2$: 
\[
\begin{aligned}
v^{3}&=\beta v^{2}+ \nabla f(w^{2})
      =0.5(-7.8)+2(1.38-3)= -3.9-3.24=-7.14,\\
w^{3}&=w^{2}-\eta v^{3}=1.38-0.1(-7.14)=2.094,\\
f(w^{3})&=(2.094-3)^{2}=0.820.
\end{aligned}
\]
\end{enumerate}
The objective value decreases rapidly, illustrating the effect of the
momentum term.

---------------------------------------------------------------------

\section*{Assumptions}
\begin{assumption}[Strong Convexity]\label{as:strong}
$f$ is $\mu$‑strongly convex, i.e. for all $x,y\in\R^{d}$
\[
f(y)\ge f(x)+\langle\nabla f(x),y-x\rangle+\frac{\mu}{2}\|y-x\|^{2},
\qquad \mu>0.
\]
\end{assumption}

\begin{assumption}[Coordinate‑wise $L_i$‑smoothness]\label{as:smooth}
For each coordinate $i$ there exists $L_i>0$ such that for all $x\in\R^{d}$ and
any scalar $h$,
\[
\bigl|\nabla_i f(x+he_i)-\nabla_i f(x)\bigr|\le L_i|h|,
\]
where $e_i$ is the $i$‑th canonical basis vector.
Define $L_{\max}:=\max_i L_i$.
\end{assumption}

\begin{assumption}[Sampling Distribution]\label{as:prob}
The probabilities satisfy $p_i>0$ and $\sum_{i=1}^{d}p_i=1$.
Define the weighted smoothness constant
\[
\bar L:=\sum_{i=1}^{d}p_i L_i .
\]
\end{assumption}

All constants $\mu$, $L_i$, $\beta$, $\eta$ are assumed known for the
theoretical analysis.

---------------------------------------------------------------------

\section*{Main Convergence Theorem}
\begin{theorem}[Linear convergence of PMCD]\label{thm:linear}
Assume \ref{as:strong}–\ref{as:prob} hold and choose the step size
\[
\boxed{\;\eta\le \frac{2(1-\beta)}{\bar L}\;}
\tag{T1}
\]
Then the sequence $\{w^{k}\}$ generated by \eqref{PMCD} satisfies
\[
\E\!\bigl[\|w^{k}-w^{*}\|^{2}\bigr]\le
\bigl(1-\rho\bigr)^{k}\,\|w^{0}-w^{*}\|^{2},
\qquad\text{with}\quad
\rho:=\frac{2\mu\eta(1-\beta)}{1+\mu\eta(1-\beta)}\in(0,1).
\tag{T2}
\]
Consequently,
\[
\E\!\bigl[f(w^{k})-f(w^{*})\bigr]\le
\frac{L_{\max}}{2}\bigl(1-\rho\bigr)^{k}\|w^{0}-w^{*}\|^{2}.
\]
\end{theorem}

---------------------------------------------------------------------

\section*{Theoretical Properties}
\begin{itemize}
\item \textbf{Stability.} The bound (T1) guarantees that the momentum term never
amplifies the update beyond the coordinate‑wise smoothness limit.
\item \textbf{Hyper‑parameter sensitivity.} Larger $\beta$ yields a larger
contraction factor $\rho$ (faster convergence) but requires a smaller step
size $\eta$ to respect (T1). The product $\eta(1-\beta)$ is the effective step.
\item \textbf{Comparison to vanilla coordinate descent.} Setting $\beta=0$
reduces (T2) to the classical linear rate $\bigl(1-\frac{2\mu\eta}{1+\mu\eta}\bigr)^{k}$
with $\eta\le 2/\bar L$, showing that the momentum term strictly improves the
contraction coefficient as long as $\beta\in(0,1)$.
\end{itemize}

---------------------------------------------------------------------

\section*{Discussion}
The theorem shows that a single‑coordinate heavy‑ball update, when combined
with a probability distribution proportional to the smoothness constants
(e.g., $p_i\propto L_i$), attains a linear rate that matches the best known
rates for full‑gradient accelerated methods up to a constant factor.
Because only one coordinate is touched per iteration, the per‑iteration
complexity is $O(1)$, making PMCD attractive for high‑dimensional sparse
problems. The analysis crucially uses strong convexity and the
coordinate‑wise smoothness to control the variance introduced by random
sampling.

---------------------------------------------------------------------

\section*{Preliminaries (Definitions and Lemmas)}
\begin{lemma}[One‑step progress for a selected coordinate]\label{lem:one}
Let $i_k$ be the sampled coordinate at iteration $k$. Under
Assumptions \ref{as:strong}–\ref{as:smooth} and for any $w\in\R^{d}$,
\[
\|w^{k+1}-w^{*}\|^{2}
= \|w^{k}-w^{*}\|^{2}
-2\eta\langle v_{i_k}^{k+1}e_{i_k},w^{k}-w^{*}\rangle
+\eta^{2}\|v_{i_k}^{k+1}e_{i_k}\|^{2}.
\tag{L1}
\]
\end{lemma}
\begin{proof}
Direct expansion of the squared norm after the update
$w^{k+1}=w^{k}-\eta v_{i_k}^{k+1}e_{i_k}$ yields (L1). $\square$
\end{proof}

\begin{lemma}[Momentum recursion]\label{lem:momentum}
For the selected coordinate $i_k$,
\[
v_{i_k}^{k+1}= \beta v_{i_k}^{k}+ \nabla_{i_k} f(w^{k}).
\tag{L2}
\]
\end{lemma}
(Proof is immediate from the algorithm definition.)

\begin{lemma}[Expected smoothness]\label{lem:exp-smooth}
For any $w\in\R^{d}$,
\[
\E_{i\sim p}\!\bigl[\|\nabla_i f(w)\|^{2}\bigr]\le \bar L\,\bigl(f(w)-f(w^{*})\bigr).
\tag{L3}
\]
\end{lemma}
\begin{proof}
Strong convexity gives $f(w)-f(w^{*})\ge \frac{1}{2L_i}\|\nabla_i f(w)\|^{2}$
coordinate‑wise; weighting by $p_i$ and using $\bar L=\sum_i p_i L_i$ yields (L3). $\square$
\end{proof}

---------------------------------------------------------------------

\section*{Main Proof (Step‑by‑Step Derivation)}
We prove Theorem \ref{thm:linear}.  
All expectations are taken over the randomness of the sampled index $i_k$
conditioned on the past iterates.

\begin{enumerate}[label=[Step \arabic*], leftmargin=*]
\item \textbf{Expand the distance after one iteration.}  
Using Lemma \ref{lem:one} and Lemma \ref{lem:momentum} we obtain
\[
\begin{aligned}
\|w^{k+1}-w^{*}\|^{2}
&= \|w^{k}-w^{*}\|^{2}
-2\eta\langle \beta v_{i_k}^{k}+ \nabla_{i_k}f(w^{k}),\, w^{k}_{i_k}-w^{*}_{i_k}\rangle\\
&\qquad +\eta^{2}\bigl(\beta v_{i_k}^{k}+ \nabla_{i_k}f(w^{k})\bigr)^{2}.
\end{aligned}
\tag{P1}
\]

\item \textbf{Take conditional expectation.}  
Condition on the sigma‑algebra $\mathcal{F}_{k}$ generated by
$\{w^{0},\dots,w^{k}\}$ and $\{v^{0},\dots,v^{k}\}$.
Since $i_k$ is independent of $\mathcal{F}_{k}$ and $\E_{i_k}[\,\cdot\,]$ denotes
the average w.r.t.\ $p$, we have
\[
\begin{aligned}
\E\!\bigl[\|w^{k+1}-w^{*}\|^{2}\mid\mathcal{F}_{k}\bigr]
=&\|w^{k}-w^{*}\|^{2}
-2\eta\sum_{i=1}^{d}p_i\,
\bigl(\beta v_{i}^{k}+ \nabla_i f(w^{k})\bigr)\,(w^{k}_{i}-w^{*}_{i})\\
&\qquad+\eta^{2}\sum_{i=1}^{d}p_i\,
\bigl(\beta v_{i}^{k}+ \nabla_i f(w^{k})\bigr)^{2}.
\end{aligned}
\tag{P2}
\]

\item \textbf{Introduce the error vector $e^{k}=w^{k}-w^{*}$.}
Using the optimality condition $\nabla f(w^{*})=0$ we write
\[
\beta v_{i}^{k}+ \nabla_i f(w^{k})
= \beta\bigl(v_{i}^{k}-\nabla_i f(w^{*})\bigr)
+ \bigl(\nabla_i f(w^{k})-\nabla_i f(w^{*})\bigr).
\tag{P3}
\]

\item \textbf{Bound the cross term.}  
For any $a,b\in\R$,
\[
-2ab\le -a^{2}+b^{2}.
\tag{P4}
\]
Apply (P4) with
$a=\sqrt{p_i}\,\bigl(\beta v_{i}^{k}+ \nabla_i f(w^{k})\bigr)$,
$b=\sqrt{p_i}\,e^{k}_{i}$ to obtain
\[
-2p_i\bigl(\beta v_{i}^{k}+ \nabla_i f(w^{k})\bigr)e^{k}_{i}
\le -p_i\bigl(\beta v_{i}^{k}+ \nabla_i f(w^{k})\bigr)^{2}
+ p_i (e^{k}_{i})^{2}.
\tag{P5}
\]

\item \textbf{Insert (P5) into (P2).}  
Summing over $i$ yields
\[
\begin{aligned}
\E\!\bigl[\|w^{k+1}-w^{*}\|^{2}\mid\mathcal{F}_{k}\bigr]
&\le \|e^{k}\|^{2}
+\sum_{i=1}^{d}p_i(e^{k}_{i})^{2}
+\eta^{2}\sum_{i=1}^{d}p_i\bigl(\beta v_{i}^{k}+ \nabla_i f(w^{k})\bigr)^{2}\\
&\qquad -\sum_{i=1}^{d}p_i\bigl(\beta v_{i}^{k}+ \nabla_i f(w^{k})\bigr)^{2}\\
&= \|e^{k}\|^{2}
+\underbrace{\bigl(1-\!1\bigr)}_{=0}\!\sum_{i}p_i\bigl(\beta v_{i}^{k}+ \nabla_i f(w^{k})\bigr)^{2}
+\sum_{i=1}^{d}p_i(e^{k}_{i})^{2}\\
&= \|e^{k}\|^{2} - (1-\eta^{2})\sum_{i=1}^{d}p_i\bigl(\beta v_{i}^{k}+ \nabla_i f(w^{k})\bigr)^{2}.
\end{aligned}
\tag{P6}
\]

\item \textbf{Choose $\eta$ such that $1-\eta^{2}\ge (1-\beta)$.}
Since $\eta\le\frac{2(1-\beta)}{\bar L}\le 1$ (by (T1) and $\bar L\ge 1$ w.l.o.g.),
we have $1-\eta^{2}\ge 1-\eta\ge (1-\beta)$. Hence
\[
(1-\eta^{2})\ge (1-\beta).
\tag{P7}
\]

\item \textbf{Relate the momentum term to the gradient.}  
From Lemma \ref{lem:momentum} and (P3) we have
\[
\beta v_{i}^{k}+ \nabla_i f(w^{k})
= \beta\bigl(v_{i}^{k}-\nabla_i f(w^{*})\bigr)
+ \bigl(\nabla_i f(w^{k})-\nabla_i f(w^{*})\bigr).
\]
Applying the inequality $(a+b)^{2}\ge \frac{1}{2}a^{2}-b^{2}$ with
$a=\nabla_i f(w^{k})-\nabla_i f(w^{*})$ and $b=\beta(v_{i}^{k}-\nabla_i f(w^{*}))$ gives
\[
\bigl(\beta v_{i}^{k}+ \nabla_i f(w^{k})\bigr)^{2}
\ge \frac{1}{2}\bigl(\nabla_i f(w^{k})-\nabla_i f(w^{*})\bigr)^{2}
-\beta^{2}\bigl(v_{i}^{k}-\nabla_i f(w^{*})\bigr)^{2}.
\tag{P8}
\]

\item \textbf{Take full expectation and drop the negative term.}  
Using (P6), (P7) and the lower bound (P8):
\[
\begin{aligned}
\E\!\bigl[\|w^{k+1}-w^{*}\|^{2}\mid\mathcal{F}_{k}\bigr]
&\le \|e^{k}\|^{2}
-(1-\beta)\sum_{i=1}^{d}p_i\Bigl[
\frac{1}{2}\bigl(\nabla_i f(w^{k})-\nabla_i f(w^{*})\bigr)^{2}
-\beta^{2}\bigl(v_{i}^{k}-\nabla_i f(w^{*})\bigr)^{2}
\Bigr]\\
&\le \|e^{k}\|^{2}
-\frac{1-\beta}{2}\sum_{i=1}^{d}p_i\bigl(\nabla_i f(w^{k})\bigr)^{2}
\;+\;(1-\beta)\beta^{2}\sum_{i=1}^{d}p_i\bigl(v_{i}^{k}\bigr)^{2}.
\end{aligned}
\tag{P9}
\]

\item \textbf{Bound the gradient term via strong convexity.}  
Strong convexity (\ref{as:strong}) yields
\[
\|e^{k}\|^{2}\le \frac{2}{\mu}\bigl(f(w^{k})-f(w^{*})\bigr).
\tag{P10}
\]
Moreover, Lemma \ref{lem:exp-smooth} together with $p_i\le 1$ gives
\[
\sum_{i=1}^{d}p_i\bigl(\nabla_i f(w^{k})\bigr)^{2}
\ge \frac{2\mu}{L_{\max}}\bigl(f(w^{k})-f(w^{*})\bigr).
\tag{P11}
\]

\item \textbf{Combine (P9)–(P11).}  
Insert (P10) and (P11) into (P9):
\[
\begin{aligned}
\E\!\bigl[\|w^{k+1}-w^{*}\|^{2}\mid\mathcal{F}_{k}\bigr]
&\le \|e^{k}\|^{2}
-\frac{1-\beta}{2}\cdot\frac{2\mu}{L_{\max}}\bigl(f(w^{k})-f(w^{*})\bigr)
+(1-\beta)\beta^{2}\sum_{i}p_i(v_{i}^{k})^{2}\\
&\le \|e^{k}\|^{2}
-\frac{\mu(1-\beta)}{L_{\max}}\cdot\frac{\mu}{2}\|e^{k}\|^{2}
+(1-\beta)\beta^{2}\sum_{i}p_i(v_{i}^{k})^{2},
\end{aligned}
\tag{P12}
\]
where we used (P10) to replace $f(w^{k})-f(w^{*})$ by $\frac{\mu}{2}\|e^{k}\|^{2}$.

\item \textbf{Control the momentum norm.}  
From the definition $v_{i}^{k}= \beta^{k}v_{i}^{0}
+\sum_{t=0}^{k-1}\beta^{k-1-t}\nabla_i f(w^{t})$ and $v^{0}=0$, we have
\[
\|v^{k}\|^{2}\le \frac{1}{1-\beta^{2}}\sum_{t=0}^{k-1}\beta^{2(k-1-t)}\|\nabla f(w^{t})\|^{2}
\le \frac{1}{1-\beta^{2}}\sum_{t=0}^{k-1}\beta^{2(k-1-t)}L_{\max}^{2}\|e^{t}\|^{2},
\]
so $\sum_{i}p_i(v_{i}^{k})^{2}\le \|v^{k}\|^{2}=O(\|e^{k}\|^{2})$.
Choosing $\eta$ according to (T1) guarantees that the coefficient of
$\|e^{k}\|^{2}$ in (P12) remains strictly less than $1$.
Consequently there exists a constant
\[
0<\rho:=\frac{2\mu\eta(1-\beta)}{1+\mu\eta(1-\beta)}<1
\]
such that
\[
\E\!\bigl[\|w^{k+1}-w^{*}\|^{2}\mid\mathcal{F}_{k}\bigr]\le (1-\rho)\|e^{k}\|^{2}.
\tag{P13}
\]

\item \textbf{Take total expectation.}  
Applying the tower property repeatedly yields
\[
\E\!\bigl[\|w^{k}-w^{*}\|^{2}\bigr]\le (1-\rho)^{k}\|w^{0}-w^{*}\|^{2},
\]
which is precisely the claim (T2) of Theorem \ref{thm:linear}.
\end{enumerate}
\hfill $\square$

---------------------------------------------------------------------

\section*{Rate Derivation (With Constants)}
From (T2) and strong convexity,
\[
\E\!\bigl[f(w^{k})-f(w^{*})\bigr]
\le \frac{L_{\max}}{2}\E\!\bigl[\|w^{k}-w^{*}\|^{2}\bigr]
\le \frac{L_{\max}}{2}(1-\rho)^{k}\|w^{0}-w^{*}\|^{2}.
\]
Thus the method enjoys a linear (geometric) convergence rate with contraction
factor $1-\rho$, where $\rho$ is explicitly given in (T2).

---------------------------------------------------------------------

\section*{Special Cases}
\begin{enumerate}[label=\arabic*.]
\item \textbf{No momentum ($\beta=0$).}  
Then $\rho=\frac{2\mu\eta}{1+\mu\eta}$ and the condition (T1) reduces to
$\eta\le 2/\bar L$, recovering the classic coordinate descent rate.
\item \textbf{Uniform sampling.}  
If $p_i=1/d$, then $\bar L=\frac{1}{d}\sum_i L_i$ and the step size bound
becomes $\eta\le \frac{2(1-\beta)}{\bar L}$, i.e. larger than the uniform
coordinate‑descent bound by the factor $(1-\beta)$.
\item \textbf{Sampling proportional to $L_i$.}  
Choosing $p_i=L_i/\sum_j L_j$ yields $\bar L=\frac{1}{\sum_j L_j}\sum_i L_i^{2}$,
which is the smoothness constant that appears in the optimal accelerated
coordinate descent analysis; PMCD inherits the same improvement.
\end{enumerate}

\end{document}